\section{System design}

The GPU is used only twice: once to log model outputs and once for the final live benchmark. All RL training runs on the CPU in a simulator built from those logs (Figure~\ref{fig:pipeline}).

\begin{figure}[h]
\centering
\begin{tikzpicture}[
  node distance=0.35cm,
  box/.style={draw, rounded corners, minimum width=7.5cm, minimum height=0.7cm, align=center, font=\small},
  gpu/.style={box, fill=orange!20},
  cpu/.style={box, fill=blue!8},
  arr/.style={-{Stealth}, thick}]
\node[cpu] (a) {1. Prompts (train / test split)};
\node[gpu, below=of a] (b) {2. Logging pass (GPU, once per model pair) $\to$ \texttt{logs/*.npz}};
\node[cpu, below=of b] (c) {3. Simulated environment built from logs (CPU)};
\node[cpu, below=of c] (d) {4. Train VI, Q-learning, DQN, REINFORCE (CPU)};
\node[gpu, below=of d] (e) {5. Live benchmark: real spec-decoding loop + KV cache (GPU)};
\node[cpu, below=of e] (f) {6. Plots, Gradio demo, report};
\foreach \x/\y in {a/b,b/c,c/d,d/e,e/f} \draw[arr] (\x) -- (\y);
\end{tikzpicture}
\caption{System pipeline. Only the orange steps need the GPU.}
\label{fig:pipeline}
\end{figure}

\paragraph{Smart logging.} With greedy decoding the final text is the same for every policy, since it is always the target's greedy output. So for each prompt we:
\begin{enumerate}
  \item Generate the target's greedy answer once (8 prompts per batch; repetition penalty off, so the text is pure greedy).
  \item Run one forward pass of the draft model over the full text.
  \item Store, at every position $t$, the draft's top probability $p_t$, its entropy $H_t$, and $\text{match}_t$ = whether the draft's top token equals the true next token.
\end{enumerate}
In the simulator, a round that starts at position $t$ with $s$ drafts accepts $n$ = the number of consecutive matches from $t$, capped at $s$. Any policy can therefore be replayed at any round start without re-running the language models. Logging costs two forward passes per prompt and takes minutes.

One approximation remains: after the first wrong draft token, a real draft model would continue from its own mistake, while the log continues from the true text. Those tokens are rejected anyway, and the live benchmark checks that the gap is small.

\paragraph{Live benchmark.} A real draft$\to$verify loop with a key-value cache. Each model keeps a fixed-size cache and one recorded CUDA graph per input length (1--10 tokens), which removes the per-step kernel-launch overhead; after each verify the cache position is moved back to the accepted length, so rejected drafts are forgotten. The target-alone baseline uses the same machinery. Correctness test: the output must equal the target's own greedy output on 20 test prompts. Timing uses 3 warm-up runs, CUDA synchronisation, and 3 timed runs.

\begin{table}[h]
\centering
\caption{Models and hardware.}
\label{tab:models}
\small
\begin{tabularx}{\linewidth}{@{}lXll@{}}
\toprule
Where & Draft $\to$ Target & VRAM (fp16, est.) & Role \\
\midrule
Laptop (RTX 3050, 6\,GB) & Qwen2.5-0.5B-Instruct $\to$ Qwen2.5-1.5B-Instruct & $\sim$4.5\,GB & Main pair \\
Laptop (RTX 3050, 6\,GB) & SmolLM2-135M-Instruct $\to$ SmolLM2-1.7B-Instruct & $\sim$4\,GB & 12$\times$ size gap, demo \\
\bottomrule
\end{tabularx}
\end{table}

Draft and target in each pair share a tokenizer, which speculative decoding requires.

\paragraph{Data.} 498 prompts mixed from Alpaca (167 chat instructions), HumanEval (all 164 code problems) and GSM8K (167 maths questions), so results can be compared per task type. 20\% of each task (99 prompts) is held out as the test set; agents train on the other 399.
